% \section{Conclusion and Future Work}
% In a future where agents team with humans in high-stakes domains \citep{kim2025srt}, accelerate science and technology research \citep{gottweis2025towards}, and empower creative endeavors \citep{waikar_2021}, it is hard to imagine how an agent incapable of coordination would contribute to such a future, however strong their individual capabilities.

% Our work demonstrates that coordination, not raw coding ability, is a central bottleneck for multi-agent software development. Through benchmarking frontier language models on \dataset, we show that when agents work in teams, success rates drop substantially compared to a single agent performing the same workload. This \emph{curse of coordination} stems from models' inability to communicate effectively, honor commitments, and maintain accurate expectations about their partners---failures that persist across tasks of varying difficulty.

% Although we focus on software development in this paper, our findings provide a starting point for investigating coordination in any domain involving role and resource conflicts under partial observability. We expect that the lack of social intelligence---the ability to understand others, communicate effectively, and coordinate actions---will remain an Achilles' heel limiting the real-world deployment of agents as teammates until these capabilities are explicitly developed.


\section{Conclusion and Future Work}

In a future where agents team with humans in high-stakes domains \citep{kim2025srt}, accelerate science and technology research \citep{gottweis2025towards}, and empower creative endeavors \citep{waikar_2021}, it is hard to imagine how an agent incapable of coordination would contribute to such a future, however strong their individual capabilities.

Our work demonstrates that coordination, not raw coding ability, is a central bottleneck for multi-agent software development. Through \dataset, we show that frontier models like \texttt{GPT-5} and \texttt{Claude Sonnet 4.5} achieve only 25\% success when two agents collaborate, roughly half the success rate of a single agent performing the same workload. This \emph{curse of coordination} stems from three capability gaps: agents fail to communicate actionable information, deviate from their own commitments, and hold incorrect expectations about their partners.

Yet coordination is not beyond reach. In successful traces, we observe emergent behaviors such as role division, resource division, and negotiation that turn vague intentions into verifiable commitments. These patterns are rare but their presence suggests the underlying capability exists; the challenge is making it reliable. With multi-agent training methods, e.g. Sotopia-$\pi$ \citep{wang2024sotopia, yu2025sotopia}, we can expect these emergent behaviors to be reinforced through the success of cooperation. 

Our findings open several directions: (1) training objectives that reward coordination under partial observability, (2) lightweight protocols for verifiable commitments (e.g., shared signatures, insertion-point contracts), and (3) richer communication channels such as screen sharing to expand the modality beyond text. We release \dataset as an open benchmark to measure progress on these fronts.

Although we focus on software development, our findings generalize to any domain involving role and resource conflicts under partial observability. We expect that the lack of social intelligence, the ability to understand others, communicate effectively, and coordinate actions, will remain a fundamental barrier limiting the real-world deployment of agents as teammates until these capabilities are explicitly developed.
